% ---------------------------------------------------------------------
% 第 7 章 结论
% ---------------------------------------------------------------------
\section{结论}

本文基于三项代数 $\mathcal A_2$（TODQ）重构机械臂运动学：串联链正运动学一次 $O(n)$ 连乘同时输出位姿、twist 与加速度；误差体系定义在两项 HDQ 元素上，一次乘法同时生成右不变位姿误差与几何一致 twist 误差（定理 1），并闭合为严格级联误差运动学（定理 2）。几何一致计算力矩律使闭环满足精确耗散等式，在同一存储函数上给出三类证书：无扰时的水平集不变性与局部指数稳定（定理 3(b)）、$L_2$ 扰动下 H∞ 增益的 Schur 补判据（定理 3(c)）、$L_\infty$ 扰动下 twist 误差的均方极限界（定理 3(d)）。加速度层被证明无需误差项——期望加速度走前馈、不确定性走扰动通道，误差层仅需 DQ/HDQ 乘法。近恒等线性化 \eqref{eq:5.8} 还揭示了一个实现陷阱：旋转刚度存在 1/4 折减，不补偿则旋转/平移带宽失配。

CoppeliaSim 力矩模式仿真（§6）定量核验了上述主张：标度律 \eqref{eq:5.9} 的扰动反演在两个独立增益档上一致至 1.93\%，base 档 12.3 倍的姿态误差放大从反面确认折减补偿的必要性；本文控制律与二阶基线 \cite{Ch20} 数值等价（$\le 0.05\%$），在高速、噪声工况下优于一阶桥接基线（$+2.0\%$--$2.2\%$）；水平集余度约 2.5 个数量级，证书条件全程未被违反。与 \cite{Ch20} 的数值等价表明，本文的贡献是同性能下的严格证书与大误差几何鲁棒性，而非精度提升。

\textbf{局限与后续工作}：strictification（附录 C.4）未完成，$\Omega_c$ 前提尚未自洽闭合；稳定性结论为工作域局部（unwinding 与 $\tilde\eta=0$ 处 $A$ 的奇异为拓扑障碍）；真机实验待验证。
